% CLAUDE REPORT::WORSTCASE-DERIVATION-0001::2026-11-15T15:00-08:00
\documentclass[11pt]{article}
\usepackage{amsmath,amssymb,amsthm,booktabs}
\usepackage[margin=1in]{geometry}
\newtheorem{proposition}{Proposition}
\newtheorem{definition}{Definition}
\title{Worst-case minimum derivation cost $\overline W_R(N)$ and the conditional bridge:\\
a logical collapse audit plus exact tree-like data}
\author{Claude scheduled routine (worst-case derivation directive), run 0001}
\date{2026-11-15 (America/Los\_Angeles)}
\begin{document}
\maketitle

\noindent\textbf{Status line.} P versus NP remains \textbf{OPEN}. Nothing here promotes a theorem
toward a separation. Propositions~1--3 are PROVED/EXACT but elementary, or are KNOWN THEOREMS.
Section~4 is COMPUTATIONAL.

\section{Target, model and quantifiers}
Fix a Cook--Reckhow propositional proof system $R$: a polynomial-time verifier, sound and complete for UNSAT CNFs.
$\langle F\rangle$ is the binary encoding of $F$ and $N=|\langle F\rangle|$. $\operatorname{cost}_R(\pi)$ is the number of lines (or symbols) of $\pi$. Then
\[
\overline W_R(N)=\max_{F\in\mathrm{UNSAT},\,|\langle F\rangle|\le N}\ \min_{\pi:F\vdash_R\bot}\operatorname{cost}_R(\pi).
\]
The inner minimum takes the best derivation for free, so finding it is not charged, as the directive specifies.
The ``multisign DAG hypergraph reduction to a minimal reduced logical state'' becomes this object once each
reduction rule is fixed as an $R$-inference. A reduction sequence is then an $R$-derivation DAG ending in $\bot$.
The target bridge is
\[
\mathrm{Bridge}(R):\qquad \mathrm{SAT}\in\mathrm P\ \Longrightarrow\ \overline W_R(N)\le\operatorname{poly}(N).
\]

\paragraph{Candidate this run.} ``$\mathrm{Bridge}(R)$ is a strictly weaker sufficient step toward P$\neq$NP than a direct lower bound against all algorithms.''
\emph{Falsifiers:} (F1) $\mathrm{Bridge}(R)$ is logically equivalent to P$\neq$NP for every $R$ with a known lower bound.
(F2) A polynomial-time-solvable UNSAT family is $R$-hard, which kills pointwise transfer.
(F3) The quantifier order $\exists R\,\forall A$ versus $\forall A\,\exists R$ trivializes the bridge.

\section{Corpus reviewed; duplication audit}
\textbf{Corpus reviewed:} \texttt{AGENTS.md} (all directives, including the 2026-11-15 SAT-compression refocus),
\texttt{research/coordination/2026-11-15\_SAT\_COMPRESSION\_CIRCUITS\_FIVE\_AGENT\_REFOCUS.md},
the world-ledger notes (\texttt{research/agents/claude/2026-10-28\_*}), recent TMR commit history up to \texttt{3356b597},
and a repository search for $\overline W_R$ and ``worst-case derivation''. That search found no prior artifact for this exact max--min target.
\textbf{Duplication:} Prop.~1 restates standard proof-complexity folklore (Cook--Reckhow 1979).
The world-ledger rule applies: a statement equivalent to P$\neq$NP is a reformulation, not progress.
Historical \texttt{/341} stays FAIL as a full UNSAT authority, and \texttt{/386} stays RESTRICTED. This report does not touch either.
\textbf{What changed relative to root main:} nothing. Root \texttt{main.tex} ownership regions belong to the TMR/LSPV roles,
so this routine does not edit them. TRACKERS REVIEWED; NO MATERIAL CHANGE.
\textbf{Stale/overstated claims found:} none new.

\section{Results}
\begin{proposition}[Bridge collapse; PROVED/EXACT]
Suppose $\overline W_R(N)\ge N^{\omega(1)}$ is proved unconditionally for $R$. This holds for tree-like and general resolution (Haken 1985),
cutting planes (Pudl\'ak 1997), bounded-depth Frege (Ajtai; Beame--Impagliazzo--Kraj\'\i\v{c}ek--Pitassi--Pudl\'ak--Woods)
and Polynomial Calculus (Razborov). Then $\mathrm{Bridge}(R)\iff \mathrm P\neq\mathrm{NP}$.
\end{proposition}
\begin{proof}
If P$\neq$NP, the antecedent is false and the implication holds vacuously.
If $\mathrm{Bridge}(R)$ holds, the false consequent gives $\neg(\mathrm{SAT}\in\mathrm P)$ by contraposition.
\end{proof}
\noindent So for every system where the max--min lower bound is already known, the ``weakest sufficient bridge'' is not weaker than the goal.
It is the goal. Any proof of it is a proof of P$\neq$NP, and it inherits every barrier: relativization, natural proofs and algebrization.

\begin{proposition}[Where the bridge has content; KNOWN THEOREM]
(a) Cook--Reckhow: some $R$ has $\overline W_R\le\operatorname{poly}$ iff NP$=$coNP.
(b) Cook (1975) and Cook--Kraj\'\i\v{c}ek (2007): if PV proves that a polynomial-time algorithm decides SAT, then Extended Frege is polynomially bounded.
So $\mathrm{Bridge}(\mathrm{EF})$ holds for every SAT algorithm whose correctness is PV-provable.
\end{proposition}
\noindent Consequence: the bridge has non-vacuous content only for strong systems such as EF or stronger, where a superpolynomial
$\overline W_R$ lower bound is itself a famous open problem. A 2025--26 literature check found no such lower bound.
Lu--Santhanam--Tzameret (ITCS~2026) get AC$^0[p]$-Frege lower bounds only for specific families, infinitely often.
Res($\oplus$) has superpolynomial size lower bounds only in restricted depth regimes.
The two halves therefore trade off: weak $R$ has the lower bound but a bridge equivalent to P$\neq$NP; strong $R$ has the bridge (under PV-provability) but no lower bound.
The precise residual gap for EF is: (i) a superpolynomial EF lower bound, plus (ii) extending the bridge to SAT algorithms whose correctness is not PV-provable.

\begin{proposition}[Quantifier order; PROVED/EXACT]
For every polynomial-time $A$ deciding SAT, the system $R_A$ (empty proof; the verifier runs $A$) satisfies $\overline W_{R_A}(N)=O(1)$.
Hence $\forall A\,\exists R$ is trivial. The bridge must be stated for a fixed $R$ chosen before $A$.
A lower bound for every $R$ is $\mathrm{NP}\neq\mathrm{coNP}$, which is stronger than P$\neq$NP.
\end{proposition}

\paragraph{Legitimate scope of the user's minimization argument (KNOWN THEOREM).}
Taking the optimal sequence for free is sound for lower bounds against algorithms whose execution traces \emph{are} $R$-derivations.
DPLL traces are tree-like resolution. CDCL with restarts is polynomially equivalent to resolution (Pipatsrisawat--Darwiche; Atserias--Fichte--Thurley).
For these algorithms, the exponential lower bounds are already unconditional. Extending them to arbitrary algorithms is exactly $\mathrm{Bridge}(R)$.

\section{Computational falsifier for pointwise transfer}
Code: \texttt{worst-case-derivation/wcd\_treeres.py}; raw data: \texttt{worst-case-derivation/2026-11-15\_wcd\_treeres\_results.json}; base commit \texttt{3356b597}.
$S_{\mathrm{tree}}$ is the exact minimum number of leaves of a DPLL tree, equal to the axiom occurrences in a minimum tree-like refutation.
It is the free minimum over all branching sequences. Unit propagation is not free. Each call takes under 10\,s; there was no chunking beyond one call per instance.
Tseitin instances use the 3-regular circulant graph $C_n$ plus diameters, with charge 1 at vertex 0 (so they are UNSAT).
The Gaussian-elimination cost counts F$_2$ row XORs on $(n_{\mathrm{vars}}+1)$-bit rows.
\begin{center}\begin{tabular}{lrrrrr}\toprule
instance & vars & clauses & $S_{\mathrm{tree}}$ & Gauss row-ops & Gauss bit-ops\\\midrule
PHP$^3_2$ & 6 & 9 & 11 & -- & --\\
PHP$^4_3$ & 12 & 22 & 43 & -- & --\\
Tseitin $K_4$ & 6 & 16 & 28 & 5 & 35\\
Tseitin circ(6) & 9 & 24 & 76 & 11 & 110\\
Tseitin circ(8) & 12 & 32 & 128 & 19 & 247\\
rand3 $n{=}10$, $m{=}60$, seed 1 & 10 & 60 & 13 & -- & --\\
\bottomrule\end{tabular}\end{center}
Two random seeds ($n{=}8$ seed 1, $n{=}10$ seed 2) were satisfiable and are excluded.
COMPUTATIONAL reading: on Tseitin, the free-optimal tree-like cost grows much faster than the polynomial cost of an F$_2$ algorithm.
Asymptotically, Tseitin on expanders needs $2^{\Omega(n)}$ resolution size (Urquhart 1987), yet it is decidable in polynomial time by linear algebra.
PHP is resolution-hard (Haken) but has polynomial cutting-planes proofs. So ``$R$-hard family $\Rightarrow$ algorithm-hard family'' is FALSIFIED pointwise (F2).
This is the proof-complexity analogue of the DISJ control. Small-$n$ numbers are not asymptotic evidence.

\section{Four-gate table}
\begin{tabular}{p{3.2cm}p{11cm}}\toprule
DISJ sanity & PASS for Props.~1--3, which make no rank-to-hardness claim. FAIL for any pointwise ``$R$-hard $\Rightarrow$ hard'' transfer: Tseitin and PHP are tractable controls.\\
Relativization & Props.~1 and 3 are pure logic and relativize. A proof of $\mathrm{Bridge}(R)$ for a lower-bounded $R$ is a proof of P$\neq$NP, so it must be nonrelativizing (Baker--Gill--Solovay). OPEN where such a step would come from.\\
Simultaneous DTISP & NOT APPLICABLE: no FLMV/time--space contradiction is claimed.\\
Configuration counting & NOT APPLICABLE as a lower bound. Remark: a DPLL run's time is at least its tree size, which is consistent and gives only an upper-bound direction for machines.\\
\bottomrule\end{tabular}

\section{Resource ledger}
The measure is lines, i.e.\ leaves and resolution steps. Strategy discovery is uncharged by design. The F$_2$ elimination counts are charged at bit level.
Peak memo states (DP memory) are in the raw JSON; the maximum was 400\,273 states, for Tseitin circ(8).

\section{Adopted / rejected / retired}
Adopted: the max--min $\overline W_R$ as the correct formalization of the user's ``minimal work over all valid reduction sequences'', which is a sound lower-bound measure for $R$-trace algorithms.
Rejected: treating $\mathrm{Bridge}(R)$ for resolution, cutting planes, bounded-depth Frege or PC as a ``weaker'' sufficient step (Prop.~1).
Retired: pointwise hard-family transfer (Section~4).

\section{Next lemma (OPEN)}
Find the weakest $R$ for which a \emph{constructive} bridge is provable without PV-provability of correctness.
Concretely: if SAT$\in$P via an algorithm $A$, do the CNFs ``$A$ is correct on inputs of length $N$'' have polynomial-size EF proofs?
This is the uniform-reflection gap. Alternatively, find a family on which every known polynomial-time SAT-algorithm paradigm induces polynomial $R$-proofs.
Natural shape: an ``algorithm $\to$ proof'' simulation for the algebraic $+$ resolution hybrids (Res($\oplus$), CP) that would be needed to absorb Gaussian elimination and counting.

\paragraph{Sources.} Cook--Reckhow 1979; Cook 1975; Cook--Kraj\'\i\v{c}ek 2007 (via arXiv:1905.12935); Haken 1985; Urquhart 1987; Pudl\'ak 1997;
Lu--Santhanam--Tzameret, ECCC TR25-134 / ITCS 2026; Bhattacharya--Chattopadhyay arXiv:2507.23008; Baker--Gill--Solovay 1975.

\noindent P versus NP remains OPEN.
\end{document}
% END_OF_VERIFIED_REPORT::CLAUDE::WORSTCASE-DERIVATION-0001
